\documentclass[11pt, twocolumn, landscape, a4paper]{article}

\usepackage[margin=1cm]{geometry}
\usepackage[utf8]{inputenc}
\usepackage[spanish]{babel}
\usepackage{amssymb, amsmath, amsbsy}
\usepackage{mathpazo}
\usepackage[pdftex]{hyperref}
\hypersetup{pdftitle={Ondulatoria},pdfauthor={Luis Carrillo Gutiérrez}}

\renewcommand{\familydefault}{\sfdefault}
\pagestyle{empty}

\begin{document}
\subsection*{Ondulatoria}

% Movimiento harmónico simples
\subsubsection*{Movimiento armónico simple}

$$
x = A \cdot \cos (\omega t + \Theta_{0})
$$

$$
v = - \omega \cdot A \cdot \sen (\omega t + \Theta_{0})
$$

$$
a = -\omega^{2}\,x
$$

% Velocidade angular de um sistema massa mola
Velocidad angular de un sistema de masa
$$
\omega = \sqrt{k / m}
$$

% Velocidade angular de um pendulo
Velocidad angular de un péndulo
$$
\omega = \sqrt{q / L}
$$

% Velocidade das ondas
\subsubsection*{Velocidad de onda}
$$
v = \dfrac{\lambda}{T}
$$

% En uma corda
En una cuerda
$$
v = \sqrt{F / \mu_{1}}
$$

% No ar (Som)
No 
$$
v = \sqrt{B / \mu}
$$

% Acustica
\subsubsection*{Acústica}

% Intensidade sonora
Intensidad sonora
$$
i = \dfrac{\Delta E}{S \cdot \Delta t} = \dfrac{P}{S}
$$

% Nivel sonoro
Nivel sónoro
$$
\beta = 10 \cdot \log \dfrac{i}{i_{0}}
$$

% Cordas e tubos sonoros
\subsubsection*{Cuerdas y tubos sonoros}

% Frequéncia de uma corda ou tubo sonoro
Frecuencia de una cuerda o tubo sonoro
$$
f = n \cdot f_{n}
$$

% Corda ou tubo aberto
Cuerda o tubo sonoro abierto
$$
f_{n} = \dfrac{v}{2L}
$$

% tubo sonoro fechado
tubo sonoro dirigido
$$
f_{n} = \dfrac{v}{4L}
$$

% Efeito Doppler
\subsubsection*{Efecto Dopler}
$$
f = f_{o} \Big(\dfrac{v \pm v_{r}}{v \pm v_{t}}\Big)
$$

% Experipencia de Young
\subsubsection*{Experiencia de Young}
$$
\lambda = dl / D
$$

\end{document}